\section{Discussion}
\label{sec:discussion}
\subsection{Separating auxiliary supervision from forecast readout}
Across the three retained seeds, auxiliary supervision increased the direct-readout mean Macro AUROC and AP but slightly decreased its mean F1. Forecast attention had higher mean AUROC, AP, and F1 than both direct controls. However, the auxiliary-supervision AUROC comparison was negative at seed 2026, and forecast attention had lower F1 than auxiliary-supervised direct readout at seed 44. These individual results remain visible alongside the summaries. Favorable average ranking does not imply improvement at every operating point or in every training realization.

None of the nine principal AUROC/AP/F1 comparisons passed within-family Holm adjustment. The full-versus-\nojepa{} Macro AUROC and AP comparisons had adjusted values of 0.18963 and 0.22488, respectively; both all-90 values were 1.0000. The conditional analysis therefore supports a descriptive ranking among these fitted predictors, rather than a multiplicity-adjusted claim of superiority. It also does not remove retrospective selection bias. The supplementary mean-pooling control had slightly lower mean Macro AUROC and AP, while its mean F1 exceeded attention by only 0.00032. All three paired metric intervals included zero. These results establish neither a necessary role for attention nor equivalence between the readout choices.

The \direct{}--\nojepa{} contrast most closely isolates future-state supervision. By comparison, the full-versus-direct contrast changes forecast input features, pooling, and classifier width together. Classifier-width-matched direct readout, stopped-gradient forecast features, copy-state readout, and independent holdout forecast evaluation would help distinguish these contributions. The main models predict embeddings rather than generate MRI for classification. The internal name \nojepa{} denotes removal of the future-state objective while retaining the shared fusion network and transition queries.

\subsection{What additional visits change}
The full model's mean AUROC increased modestly across T0--T2, but its paired T2--T0 interval included zero. Its mean AP decreased slightly, while F1 decreased at T1 before increasing at T2. Changes in true- and false-positive predictions explain why added observations, ranking, and a fixed-threshold decision rule need not improve together. Outcome-conditioned histograms summarize probability changes within each seed-specific ensemble; they do not demonstrate that every patient's prediction improves or that observed score changes are treatment effects.

T0 imaging preceded randomization, while the model is conditioned on the eventual assigned arm. This remains a retrospective conditional scenario. The arm encoding does not capture delivered dose, cycle timing, dose reductions, adherence, or treatment changes. Altering the arm token cannot identify an individual treatment effect without suitable causal assumptions and exposure data. Complete T0--T2 holdout eligibility also restricts the population to which the longitudinal comparison applies.

\subsection{Classification, probability quality, and hypothetical utility}
The additional metrics expose trade-offs behind the ranking means. Compared with auxiliary-supervised direct readout, forecast attention had higher mean sensitivity, NPV, balanced accuracy, and MCC but lower specificity, PPV, and accuracy. The specificity decrease passed both the 81-test secondary-family correction (0.02592) and the all-90 correction (0.02880), whereas the sensitivity increase did not pass its secondary-family criterion. This adverse operating-point difference should accompany the favorable discrimination summary.

Forecast attention also had worse NLL, Brier score, and ECE point estimates than both direct-readout controls, and the mean-pooling control also had better values. Its higher ranking means therefore do not imply better predicted probabilities. The OOF Youden analysis slightly increased F1 for \nojepa{} but decreased it for the other variants. Development OOF scores originate from one held-out fold model per patient, whereas test scores average five models within a seed. Transferring thresholds between these score distributions is a sensitivity analysis rather than independent ensemble calibration.

Net-benefit orderings changed with the hypothetical threshold: direct readout was higher at 0.30 and 0.50, while forecast attention was higher at 0.70. These are seed means under a generic positive-action utility, not evaluations of a defined care pathway. A clinically motivated threshold requires consequences for false-positive and false-negative actions to be specified and calibration to be fitted without holdout leakage.

\subsection{Seed variation, subgroups, and forecasting diagnostics}
The mean and sample standard deviation across three seeds describe three training realizations using identical development folds and holdout patients. They do not turn seeds into independent clinical replications. Patient-bootstrap intervals and tests answer a different question: how the mean metric of these fixed predictors varies under patient resampling. Because seeds are not resampled and models are not refitted, those intervals do not estimate uncertainty over the training algorithm or future seed choices.

Subgroup summaries are exploratory, use small denominators, and do not test subtype interactions. All four HR/HER2 full-versus-control Macro AUROC intervals included zero. The positive point estimates do not establish subtype-specific benefit; low sensitivity at 0.5 in the HR-positive/HER2-negative subgroup also persists. Source namespaces share one clinical release and do not provide external validation. Crop analysis is similarly constrained: the 11 patients with partial clipping contain only one positive outcome, preventing reliable discrimination assessment under the minimum-outcome rule. Reporting the denominator and suppressing AUROC/AP is preferable to interpreting an unstable estimate.

Archived fixed-training-probe diagnostics are separate from the pCR holdout analysis. Their copy-normalized skills and gradient norms are calculated in each snapshot's own EMA representation, averaged over folds within seed, and then summarized across seeds. They do not ensemble latent states or establish anatomical forecast fidelity. At epoch 20, auxiliary-supervised direct models had higher mean copy-normalized skill than forecast attention at all three horizons, despite the latter's higher mean pCR AUROC. This weakens an explanation based solely on more faithful forecasts. Gradient routing can support an implementation-level mechanism, but gradient magnitude is not a biological effect size. Independent forecast evaluation at the selected pCR checkpoints remains necessary.

\subsection{Limitations and a confirmatory study}
\label{sec:limitations}
The seed set 44, 2026, and 3407 was obtained by excluding seed 42 after holdout results had been available. This reporting-policy change, within-seed probability ensembling, and the statistical procedures are retrospective. Patient-bootstrap inference conditions on these fitted ensembles and fixed OOF thresholds. It excludes retraining, seed-set selection, aggregation choice, model selection, and threshold-selection uncertainty. Three selected seeds do not establish broad training stability, and their sample standard deviations have limited precision. The reported mean describes three five-model predictors; it is not the performance of one fifteen-model ensemble or a model trained on the full development cohort.

The 95\% patient intervals are unadjusted. Holm corrections for nine principal and 81 secondary tests, together with the all-90 sensitivity analysis, do not undo retrospective choices. Null-centered bootstrap tests are approximate, with finite-sample limitations particularly for non-smooth statistics such as AP, F1, and binned ECE. The absence of adjusted principal-model differences and the adverse specificity result should both remain visible. Nonsignificance does not establish equivalence. A confirmatory study should lock seeds, aggregation, model and threshold selection, and statistical comparisons before evaluating independent patients.

The holdout is an internal downstream split. Twenty-nine patients had belonged to an earlier development split. Current supervised models were retrained using the new patient-disjoint folds and were not warm-started from those earlier pCR models, but codec and localization-model pretraining have not been fully audited for current-holdout overlap. This prevents a claim that every pipeline stage was evaluated on wholly unseen patients.

Fixed baseline localization and cropping may omit multifocal disease or changes outside the field of view. Geometry and array checks do not establish anatomical registration accuracy. Holdout patients have complete T0--T2 imaging, whereas development retains incomplete histories; missingness, acquisition variation, and subtype prevalence may limit transportability. Archived training probes do not replace best-checkpoint holdout forecast evaluation and have small long-horizon denominators.

The study remains an internal architectural comparison. Matched clinical-only, conventional imaging, and independent longitudinal baselines are still needed, along with external validation, anatomical quality review, and independently assessed calibration. Precise pCR pathology criteria, scan intervals, ethics determination, and upstream provenance also require final documentation before submission.
